% ================================================================================
% Paquetes y comandos definidos por el usuario
% ================================================================================

% Permite definir tipos de columna personalizados en tablas
\usepackage{array}

% Reglas horizontales de calidad tipográfica para cuadros en formato científico
% (sin líneas verticales ni color), conforme a la guía de la universidad.
\usepackage{booktabs}

% La guía de la universidad denomina "cuadro" a las tablas, y así se titula el
% listado correspondiente. Se ajusta el rótulo de los pies para que coincida con
% el listado y con las referencias cruzadas del texto.
\addto\captionsspanish{\renewcommand{\tablename}{Cuadro}}

% Columna de párrafo alineada a la izquierda (sin justificar). Se usa en cuadros
% con columnas angostas, donde justificar el texto produce espacios excesivos.
\newcolumntype{L}[1]{>{\raggedright\arraybackslash}p{#1}}

% ================================================================================
% NUMERACIÓN Y SANGRÍA DE FIGURAS Y CUADROS
% ================================================================================
% La clase 'report' numera figuras y cuadros por capítulo (8.1, 8.2, ...) y los
% reinicia en cada uno. Se desvincula el contador del capítulo para que la
% numeración sea corrida a lo largo de todo el trabajo (1, 2, 3, ...).
\counterwithout{figure}{chapter}
\counterwithout{table}{chapter}

% Los listados de figuras y cuadros heredan la sangría de un nivel interior. Se
% redefinen sus entradas para que queden alineadas al margen izquierdo, como
% corresponde a un elemento de primer nivel.
\makeatletter
\renewcommand*\l@figure{\@dottedtocline{1}{0em}{2.3em}}
\renewcommand*\l@table{\@dottedtocline{1}{0em}{2.3em}}
\makeatother
